% GENERATED FILE. Edit canonical lessons, then run npm run generate:books.
% canonical-source-hash: fnv1a64:e844ed96fe7579b0
% canonical-lessons: TE-C501-santosapettu, TE-C501-ascaryaparacu, TE-C501-satkarincu, TE-C501-meppincu, TE-C501-utsahaparacu, TE-R501-review-494-501

\chapter{Making Others Feel}
\label{ch:te-making-others-feel}
\hlchaptermodality{501}

\begin{chapteropening}
\textbf{By the end of this chapter:} \emph{I can say how someone pleased, surprised, honoured, impressed or cheered others.}

Complete the last lesson of chapter 501: I can say how someone pleased, surprised, honoured, impressed or cheered others.
\end{chapteropening}

\section[\texorpdfstring{\te{సంతోషపెట్టు} (santōṣapeṭṭu)}{సంతోషపెట్టు (santōṣapeṭṭu)}]{\te{సంతోషపెట్టు} (santōṣapeṭṭu) — to make happy; to please}

\label{lesson:TE-C501-santosapettu}



\begin{quote}
Before the new one: say the Telugu for to thrill, then the Telugu for cheerfulness.
\end{quote}

\subsection*{You'll want to know: \te{సంతోషపెట్టు}}
\textbf{\te{సంతోషపెట్టు}} — \emph{santōṣapeṭṭu} — \textquotedblleft{}to make happy; to please\textquotedblright{}.

\begin{quote}
\te{పుట్టినరోజున} \te{అన్నయ్య} \te{అమ్మను} \te{సంతోషపెట్టాడు}.
\end{quote}

\begin{quote}
\emph{puṭṭinarōjuna annayya ammanu santōṣapeṭṭāḍu.}
\end{quote}

\begin{quote}
On her birthday my brother made mum happy.
\end{quote}

Past (I …): \te{సంతోషపెట్టాను}, \emph{santōṣapeṭṭānu}.

\te{సంతోషం}, happiness, + \te{పెట్టు}, to put: to make someone happy.

\begin{cousinweb}
\textbf{\te{సంతోషపెట్టు}} (\emph{santōṣapeṭṭu}) joins \te{సంతోషం}, from Sanskrit \emph{saṁtōṣa}, \textquotedblleft{}satisfaction\textquotedblright{} (CDIAL 12895), to the native verb \emph{peṭṭu}, \textquotedblleft{}to put\textquotedblright{}.
\end{cousinweb}

\begin{grammarlens}[title={What you've built}]
One more word for describing what you lived and felt.
\end{grammarlens}

\subsection*{Your turn}
\begin{itemize}
\raggedright
  \item \emph{Say it:} \emph{santōṣapeṭṭu}
  \item \emph{Say it:} \emph{santōṣapeṭṭu}, once more
  \item \emph{Say it:} \emph{puṭṭinarōjuna annayya ammanu santōṣapeṭṭāḍu.} — On her birthday my brother made mum happy.
  \item \emph{Recall:} say the Telugu for to thrill, then the Telugu for cheerfulness, then say \emph{santōṣapeṭṭu} again
\end{itemize}

\begin{tcolorbox}[breakable,colback=teal!4,colframe=teal!35!black,title={Before you move on}]
What does \te{సంతోషపెట్టు} mean? (\textquotedblleft{}To make happy; to please\textquotedblright{}.) Say it once more.
\end{tcolorbox}
\section[\texorpdfstring{\te{ఆశ్చర్యపరచు} (āścaryaparacu)}{ఆశ్చర్యపరచు (āścaryaparacu)}]{\te{ఆశ్చర్యపరచు} (āścaryaparacu) — to surprise; to astonish}

\label{lesson:TE-C501-ascaryaparacu}



\begin{quote}
Before the new one: say the Telugu for cheerfulness, then the Telugu for to make happy.
\end{quote}

\subsection*{You'll want to know: \te{ఆశ్చర్యపరచు}}
\textbf{\te{ఆశ్చర్యపరచు}} — \emph{āścaryaparacu} — \textquotedblleft{}to surprise; to astonish\textquotedblright{}.

\begin{quote}
\te{ఆమె} \te{పాట} \te{అందరినీ} \te{ఆశ్చర్యపరిచింది}.
\end{quote}

\begin{quote}
\emph{āme pāṭa andarinī āścaryaparicindi.}
\end{quote}

\begin{quote}
Her song astonished everyone.
\end{quote}

Past (I …): \te{ఆశ్చర్యపరిచాను}, \emph{āścaryaparicānu}.

-\te{పరచు} makes a causative from a noun: \te{ఆశ్చర్యం}, surprise; \te{ఆశ్చర్యపరచు}, to surprise someone.

\begin{cousinweb}
\textbf{\te{ఆశ్చర్యపరచు}} (\emph{āścaryaparacu}) is built from \te{ఆశ్చర్యం}, from Sanskrit \emph{āścarya}, \textquotedblleft{}wonderful\textquotedblright{} (CDIAL 1464) and \emph{paracu}, which DEDR 3853 lists as \textquotedblleft{}to cause to feel\textquotedblright{}, the causative partner of \emph{paḍu}.
\end{cousinweb}

\begin{grammarlens}[title={What you've built}]
One more word for describing what you lived and felt.
\end{grammarlens}

\subsection*{Your turn}
\begin{itemize}
\raggedright
  \item \emph{Say it:} \emph{āścaryaparacu}
  \item \emph{Say it:} \emph{āścaryaparacu}, once more
  \item \emph{Say it:} \emph{āme pāṭa andarinī āścaryaparicindi.} — Her song astonished everyone.
  \item \emph{Recall:} say the Telugu for cheerfulness, then the Telugu for to make happy, then say \emph{āścaryaparacu} again
\end{itemize}

\begin{tcolorbox}[breakable,colback=teal!4,colframe=teal!35!black,title={Before you move on}]
What does \te{ఆశ్చర్యపరచు} mean? (\textquotedblleft{}To surprise; to astonish\textquotedblright{}.) Say it once more.
\end{tcolorbox}
\section[\texorpdfstring{\te{సత్కరించు} (satkariñcu)}{సత్కరించు (satkariñcu)}]{\te{సత్కరించు} (satkariñcu) — to honour formally; to felicitate}

\label{lesson:TE-C501-satkarincu}



\begin{quote}
Before the new one: say the Telugu for to make happy, then the Telugu for to surprise.
\end{quote}

\subsection*{You'll want to know: \te{సత్కరించు}}
\textbf{\te{సత్కరించు}} — \emph{satkariñcu} — \textquotedblleft{}to honour formally; to felicitate\textquotedblright{}.

\begin{quote}
\te{బడిలో} \te{మంచి} \te{మార్కులు} \te{వచ్చిన} \te{విద్యార్థులను} \te{సత్కరించారు}.
\end{quote}

\begin{quote}
\emph{baḍilō mañci mārkulu vaccina vidyārthulanu satkariñcāru.}
\end{quote}

\begin{quote}
The school honoured the students who got good marks.
\end{quote}

Past (I …): \te{సత్కరించాను}, \emph{satkariñcānu}.

Honouring someone in public, often with a shawl and a garland. The noun is \te{సత్కారం}.

\begin{cousinweb}
\textbf{\te{సత్కరించు}} (\emph{satkariñcu}) is built on Sanskrit \emph{satkāra}, \textquotedblleft{}kind treatment, honour, hospitality\textquotedblright{} (Monier-Williams), from \emph{sat}, \textquotedblleft{}good\textquotedblright{}, and \emph{kāra}, \textquotedblleft{}doing\textquotedblright{}, with the Telugu verb ending -\te{ఇంచు}.
\end{cousinweb}

\begin{grammarlens}[title={What you've built}]
One more word for describing what you lived and felt.
\end{grammarlens}

\subsection*{Your turn}
\begin{itemize}
\raggedright
  \item \emph{Say it:} \emph{satkariñcu}
  \item \emph{Say it:} \emph{satkariñcu}, once more
  \item \emph{Say it:} \emph{baḍilō mañci mārkulu vaccina vidyārthulanu satkariñcāru.} — The school honoured the students who got good marks.
  \item \emph{Recall:} say the Telugu for to make happy, then the Telugu for to surprise, then say \emph{satkariñcu} again
\end{itemize}

\begin{tcolorbox}[breakable,colback=teal!4,colframe=teal!35!black,title={Before you move on}]
What does \te{సత్కరించు} mean? (\textquotedblleft{}To honour formally; to felicitate\textquotedblright{}.) Say it once more.
\end{tcolorbox}
\section[\texorpdfstring{\te{మెప్పించు} (meppiñcu)}{మెప్పించు (meppiñcu)}]{\te{మెప్పించు} (meppiñcu) — to impress; to win approval}

\label{lesson:TE-C501-meppincu}



\begin{quote}
Before the new one: say the Telugu for to surprise, then the Telugu for to honour formally.
\end{quote}

\subsection*{You'll want to know: \te{మెప్పించు}}
\textbf{\te{మెప్పించు}} — \emph{meppiñcu} — \textquotedblleft{}to impress; to win approval\textquotedblright{}.

\begin{quote}
\te{ఆ} \te{పిల్లవాడు} \te{తన} \te{పాటతో} \te{అందరినీ} \te{మెప్పించాడు}.
\end{quote}

\begin{quote}
\emph{ā pillavāḍu tana pāṭatō andarinī meppiñcāḍu.}
\end{quote}

\begin{quote}
That boy impressed everyone with his song.
\end{quote}

Past (I …): \te{మెప్పించాను}, \emph{meppiñcānu}.

To do something so well that others praise you; \te{మెప్పు} is the praise you earn.

\begin{cousinweb}
\textbf{\te{మెప్పించు}} (\emph{meppiñcu}) is native Dravidian: DEDR 4722 lists Telugu \emph{meppincu}, \textquotedblleft{}to please, to satisfy, to win the approval or admiration of\textquotedblright{}, beside \emph{meccu}, \textquotedblleft{}to praise\textquotedblright{}, and \emph{meppu}, \textquotedblleft{}praise\textquotedblright{}.
\end{cousinweb}

\begin{grammarlens}[title={What you've built}]
One more word for describing what you lived and felt.
\end{grammarlens}

\subsection*{Your turn}
\begin{itemize}
\raggedright
  \item \emph{Say it:} \emph{meppiñcu}
  \item \emph{Say it:} \emph{meppiñcu}, once more
  \item \emph{Say it:} \emph{ā pillavāḍu tana pāṭatō andarinī meppiñcāḍu.} — That boy impressed everyone with his song.
  \item \emph{Recall:} say the Telugu for to surprise, then the Telugu for to honour formally, then say \emph{meppiñcu} again
\end{itemize}

\begin{tcolorbox}[breakable,colback=teal!4,colframe=teal!35!black,title={Before you move on}]
What does \te{మెప్పించు} mean? (\textquotedblleft{}To impress; to win approval\textquotedblright{}.) Say it once more.
\end{tcolorbox}
\section[\texorpdfstring{\te{ఉత్సాహపరచు} (utsāhaparacu)}{ఉత్సాహపరచు (utsāhaparacu)}]{\te{ఉత్సాహపరచు} (utsāhaparacu) — to cheer on; to enthuse}

\label{lesson:TE-C501-utsahaparacu}



\begin{quote}
Before the new one: say the Telugu for to honour formally, then the Telugu for to impress.
\end{quote}

\subsection*{You'll want to know: \te{ఉత్సాహపరచు}}
\textbf{\te{ఉత్సాహపరచు}} — \emph{utsāhaparacu} — \textquotedblleft{}to cheer on; to enthuse\textquotedblright{}.

\begin{quote}
\te{అమ్మ} \te{మాటలు} \te{నన్ను} \te{ఉత్సాహపరిచాయి}.
\end{quote}

\begin{quote}
\emph{amma māṭalu nannu utsāhaparicāyi.}
\end{quote}

\begin{quote}
Mum's words cheered me on.
\end{quote}

Past (I …): \te{ఉత్సాహపరిచాను}, \emph{utsāhaparicānu}.

\te{ఉత్సాహం}, enthusiasm, + -\te{పరచు}, to make someone feel it.

\begin{cousinweb}
\textbf{\te{ఉత్సాహపరచు}} (\emph{utsāhaparacu}) is built from \te{ఉత్సాహం}, from Sanskrit \emph{utsāha}, \textquotedblleft{}power, firmness; joy\textquotedblright{} (CDIAL 1882), and \emph{paracu}, which DEDR 3853 lists as \textquotedblleft{}to cause to feel\textquotedblright{}.
\end{cousinweb}

\begin{grammarlens}[title={What you've built}]
One more word for describing what you lived and felt.
\end{grammarlens}

\subsection*{Your turn}
\begin{itemize}
\raggedright
  \item \emph{Say it:} \emph{utsāhaparacu}
  \item \emph{Say it:} \emph{utsāhaparacu}, once more
  \item \emph{Say it:} \emph{amma māṭalu nannu utsāhaparicāyi.} — Mum's words cheered me on.
  \item \emph{Recall:} say the Telugu for to honour formally, then the Telugu for to impress, then say \emph{utsāhaparacu} again
\end{itemize}

\begin{tcolorbox}[breakable,colback=teal!4,colframe=teal!35!black,title={Before you move on}]
What does \te{ఉత్సాహపరచు} mean? (\textquotedblleft{}To cheer on; to enthuse\textquotedblright{}.) Say it once more.
\end{tcolorbox}
\section[(dialogue)]{Feelings, first words}

\label{lesson:TE-R501-review-494-501}



\begin{quote}
Say the Telugu for gloom and to worry.
\end{quote}

\subsection*{Your turn}
Say these aloud:

\begin{itemize}
\raggedright
  \item gloom, to worry, to be agitated, to panic, eagerness
  \item to be afraid, to frighten, to be startled, a nightmare, to be distraught
  \item to feel sad, to hurt, to be crushed, to grieve inwardly, despair
  \item to get angry, to flare up, to snap at, fury, irritation
  \item self-respect, humiliation, to be proud, to hold back politely, to hang one's head
  \item to be jealous, to hate, to loathe, a grudge, contempt
  \item to enjoy, to be delighted, to be overjoyed, to thrill, cheerfulness
  \item to make happy, to surprise, to honour formally, to impress, to cheer on
\end{itemize}

\begin{tcolorbox}[breakable,colback=teal!4,colframe=teal!35!black,title={Before you move on}]
Gloom? (\textbf{\te{దిగులు}}, \emph{digulu}.) To worry? (\textbf{\te{చింతించు}}, \emph{cintiñcu}.) To be agitated? (\textbf{\te{కలవరపడు}}, \emph{kalavarapaḍu}.) To panic? (\textbf{\te{కంగారుపడు}}, \emph{kaṅgārupaḍu}.) Eagerness? (\textbf{\te{ఆత్రుత}}, \emph{ātruta}.) To be afraid? (\textbf{\te{భయపడు}}, \emph{bhayapaḍu}.) To frighten? (\textbf{\te{భయపెట్టు}}, \emph{bhayapeṭṭu}.) To be startled? (\textbf{\te{ఉలిక్కిపడు}}, \emph{ulikkipaḍu}.) A nightmare? (\textbf{\te{పీడకల}}, \emph{pīḍakala}.) To be distraught? (\textbf{\te{తల్లడిల్లు}}, \emph{tallaḍillu}.) To feel sad? (\textbf{\te{బాధపడు}}, \emph{bādhapaḍu}.) To hurt? (\textbf{\te{బాధపెట్టు}}, \emph{bādhapeṭṭu}.) To be crushed? (\textbf{\te{కుంగిపోవు}}, \emph{kuṅgipōvu}.) To grieve inwardly? (\textbf{\te{కుమిలిపోవు}}, \emph{kumilipōvu}.) Despair? (\textbf{\te{నిస్పృహ}}, \emph{nispṛha}.) To get angry? (\textbf{\te{కోప్పడు}}, \emph{kōppaḍu}.) To flare up? (\textbf{\te{మండిపడు}}, \emph{maṇḍipaḍu}.) To snap at? (\textbf{\te{కసురుకో}}, \emph{kasurukō}.) Fury? (\textbf{\te{ఆవేశం}}, \emph{āvēśaṁ}.) Irritation? (\textbf{\te{చిరాకు}}, \emph{cirāku}.) Self-respect? (\textbf{\te{ఆత్మగౌరవం}}, \emph{ātmagauravaṁ}.) Humiliation? (\textbf{\te{అవమానం}}, \emph{avamānaṁ}.) To be proud? (\textbf{\te{గర్వపడు}}, \emph{garvapaḍu}.) To hold back politely? (\textbf{\te{మొహమాటపడు}}, \emph{mohamāṭapaḍu}.) To hang one's head? (\textbf{\te{తలదించుకో}}, \emph{taladiñcukō}.) To be jealous? (\textbf{\te{అసూయపడు}}, \emph{asūyapaḍu}.) To hate? (\textbf{\te{ద్వేషించు}}, \emph{dvēṣiñcu}.) To loathe? (\textbf{\te{అసహ్యించుకో}}, \emph{asahyiñcukō}.) A grudge? (\textbf{\te{పగ}}, \emph{paga}.) Contempt? (\textbf{\te{చిన్నచూపు}}, \emph{cinnacūpu}.) To enjoy? (\textbf{\te{ఆనందించు}}, \emph{ānandiñcu}.) To be delighted? (\textbf{\te{మురిసిపోవు}}, \emph{murisipōvu}.) To be overjoyed? (\textbf{\te{పొంగిపోవు}}, \emph{poṅgipōvu}.) To thrill? (\textbf{\te{పులకించు}}, \emph{pulakiñcu}.) Cheerfulness? (\textbf{\te{ఉల్లాసం}}, \emph{ullāsaṁ}.) To make happy? (\textbf{\te{సంతోషపెట్టు}}, \emph{santōṣapeṭṭu}.) To surprise? (\textbf{\te{ఆశ్చర్యపరచు}}, \emph{āścaryaparacu}.) To honour formally? (\textbf{\te{సత్కరించు}}, \emph{satkariñcu}.) To impress? (\textbf{\te{మెప్పించు}}, \emph{meppiñcu}.) To cheer on? (\textbf{\te{ఉత్సాహపరచు}}, \emph{utsāhaparacu}.)
\end{tcolorbox}
